% !PRINTED 509
\weberpage{509}

%<<CONT
only \gesperrt{proper} solutions of (8) come into consideration
(\S\,114). But by (9) every proper solution leads to a
system of values $\alpha, \beta$, and, through the congruences

\weq{10}{c\alpha \equiv Cx, \qquad c\beta \equiv \dfrac{Bx-y}{2} \quad (\mathrm{mod}\ n)}

to a determinate value of $c\ (\mathrm{mod}\ n)$, which may in addition be
assumed divisible by 8. If conversely $x, y$
satisfy these conditions, then $\alpha$ and $\beta$ have no common divisor, as appears from

\[ \alpha\,\dfrac{Bx-y}{2} - \beta Cx = -n \]

For according to this, an odd divisor of $\alpha$
and $\beta$ would have to divide $n$, and would therefore, since $A$ was
assumed relatively prime to $n$, divide $x$ and consequently $y$.

It follows further, if we leave aside the two exceptional cases
$D=-3$, $D=-4$, which do not concern us here at
all, because for these $j(\omega)$ is rational, that \gesperrt{different} solutions of (8) also lead to different values $c$.
For if we assume that one and the same value of $c$
could lead to two different systems $(\alpha, \beta, \gamma, \delta)$, then
from (5) a relation of the form

\[ \omega = \dfrac{\gamma+\delta\omega}{\alpha+\beta\omega} \]

would follow, and this substitution, if it is not the identical one,
is possible only for the two exceptional cases mentioned.

\gesperrt{Accordingly the number of the $v$-values in (2) which by
(3) become equal to $u$ is as great as the number of the proper solutions of the equation (8).}

If the equation (8) has only \gesperrt{one} such solution, the
corresponding value of $M$ is rationally expressible by $u$. But if
it has several solutions, the corresponding values of $M$
are the roots of a rational equation:

\weq{11}{\Phi(M,u) = 0,}

whose degree with respect to $M$ is just as great as the number
of these solutions (\S\,134).

Furthermore $u$ satisfies the class equation

\[ H_{-D}(u) = 0. \]
